\chapter{Intermediary Asset Pricing}
\label{ch:intermediary}

\begin{projectconnection}[This Is the Core Thesis of Your Project]
Every feature you engineer in Phase~2 is a proxy for the intermediary constraints described in this chapter. He and Krishnamurthy (2013) predict that dealer capital constraints drive risk premia. Adrian, Etula, and Muir (2014) show that a single intermediary-leverage factor prices the cross-section of equities, momentum, and bonds with $R^2 = 77\%$. He, Kelly, and Manela (2017) extend this result across seven asset classes---equities, government and corporate bonds, sovereign bonds, options, CDS, commodities, and FX---using a single intermediary capital ratio. Your project tests this theory with proprietary daily data from SecDB: the data the academics wished they had.
\end{projectconnection}

% ──────────────────────────────────────────────────────────────
\section{Classical vs.\ Intermediary Asset Pricing}
\label{sec:classical-vs-intermediary}
% ──────────────────────────────────────────────────────────────

\begin{prereq}[Representative Agent]
Classical macro-finance models assume a single ``representative household'' whose preferences determine all asset prices. This is a modeling device: it does not mean everyone is identical, but rather that aggregate prices behave \emph{as if} a single agent with average preferences were trading. The SDF from Chapter~\ref{ch:asset-pricing} is this agent's marginal rate of substitution between consumption today and tomorrow.
\end{prereq}

Chapter~\ref{ch:asset-pricing} developed the SDF framework: $\E_t[\SDF_{t+1} R_{i,t+1}] = 1$, where the pricing kernel $\SDF$ is driven by the marginal investor's state-dependent preferences. The classical approach assumes the marginal investor is a representative household, giving:
\begin{equation}
\SDF_{t+1} = \beta \frac{u'(C_{t+1})}{u'(C_t)}
\label{eq:consumption-sdf}
\end{equation}
\begin{itemize}[nosep]
  \item $\beta$ --- subjective discount factor (how impatient the household is; typically ${\approx}0.99$ quarterly).
  \item $u'(C_t)$ --- marginal utility of consumption at time $t$. High when consumption is low (bad times).
  \item The ratio $u'(C_{t+1})/u'(C_t)$ --- marginal rate of substitution. High when consumption falls (recession).
\end{itemize}

This is elegant, but it fails empirically. The problem is known as the \textbf{equity premium puzzle} (Mehra and Prescott, 1985).

\begin{prereq}[Risk Aversion and Utility]
A \textbf{utility function} $u(C)$ assigns a ``happiness value'' to consumption level $C$. Risk aversion means $u$ is concave: the marginal utility of an extra dollar is higher when you are poor than when you are rich. The coefficient of \textbf{relative risk aversion} (RRA) $\gamma = -C \, u''(C)/u'(C)$ measures curvature. Empirical estimates from labor supply and insurance data suggest $\gamma \approx 1$--$5$ for typical households. Higher $\gamma$ means the agent cares more about consumption fluctuations.
\end{prereq}

\begin{keyidea}[The Equity Premium Puzzle]
U.S.\ equities earned roughly 6--8\% per year more than Treasury bills over 1926--2023 (Dimson, Marsh, and Staunton, 2023). But aggregate consumption growth is extremely smooth: its annual standard deviation is only about 1--2\%. To generate a 6\% equity premium from equation~\eqref{eq:consumption-sdf}, you need the household's risk aversion $\gamma$ to exceed 50---an order of magnitude above plausible estimates. Alternatively, you need the correlation between consumption and stock returns to be much higher than it is in the data. Neither works. The household SDF simply does not move enough to explain observed risk premia.
\end{keyidea}

The intermediary approach offers a resolution. Instead of asking ``what does the representative household fear?'' it asks: \emph{who is actually the marginal investor in most markets?}

\begin{prereq}[Marginal Investor]
The \textbf{marginal investor} is the agent whose trading activity sets the price at the margin. In equity markets, this might plausibly be households (via mutual funds). But in derivatives, credit, FX, and fixed income, the marginal investor is almost always a \textbf{financial intermediary}---a broker-dealer, market maker, or bank. These institutions hold the inventory, provide liquidity, and absorb order-flow imbalances. If their constraints tighten, they pull back, and risk premia rise.
\end{prereq}

\begin{keyidea}[The Intermediary Asset Pricing Hypothesis]
If financial intermediaries---not households---are the marginal investors in most asset markets, then the relevant SDF is the intermediary's marginal value of wealth, not the household's marginal utility of consumption. The intermediary SDF depends on the intermediary's \emph{capital constraints}: its equity, leverage, VaR limits, and regulatory requirements. When capital is scarce, the intermediary's marginal value of a dollar is high, the SDF is high, and risk premia rise. This is the theoretical foundation for every feature you will build.
\end{keyidea}

\begin{intuition}[Dealers as the Fire Department]
Think of dealers as a city's fire department. In normal times, the department has ample capacity---trucks sit idle, firefighters train, response times are fast. When a major fire breaks out, the department is stretched. If a \emph{second} fire starts while the first is raging, the city has a serious problem: the marginal cost of deploying another truck is enormous. In financial markets, ``fires'' are large sell-offs or liquidity demands. In normal times, dealers absorb these easily. But when their capital is depleted (the first fire), even a modest additional shock (the second fire) causes outsized price moves. VaR utilization near 100\% is the financial equivalent of all trucks deployed. Your project measures exactly this.
\end{intuition}

The key difference between classical and intermediary pricing can be stated precisely. Classical models write $\SDF_{t+1} = g(C_{t+1}/C_t)$ for some function $g$ of aggregate consumption growth. Intermediary models write:
\begin{equation}
\SDF_{t+1} = h(w_{t+1})
\label{eq:intermediary-sdf-general}
\end{equation}
where $w_{t+1}$ is the intermediary's wealth or capital ratio, and $h$ is a function that increases sharply as $w$ falls toward the constraint boundary. Both satisfy $\E_t[\SDF_{t+1} R_{i,t+1}] = 1$ from Chapter~\ref{ch:asset-pricing}. They differ in \emph{whose} state variable drives $\SDF$.

% ──────────────────────────────────────────────────────────────
\section{He and Krishnamurthy (2013): The Theoretical Foundation}
\label{sec:he-krishnamurthy}
% ──────────────────────────────────────────────────────────────

He and Krishnamurthy (2013, AER) provide the theoretical model that your project is built on. The model has two agents: households and intermediaries.

\begin{prereq}[Balance Sheet Mechanics]
A firm's balance sheet satisfies: $\text{Assets} = \text{Equity} + \text{Debt}$. \textbf{Equity} (or ``capital'') is the firm's own money; \textbf{debt} is borrowed money. \textbf{Leverage} is the ratio $\text{Assets}/\text{Equity}$. A firm with \$100 in assets, \$10 in equity, and \$90 in debt has leverage of 10$\times$. If asset values drop by 5\% (to \$95), equity drops by 50\% (to \$5)---leverage amplifies losses. This is why bank capital matters: small asset losses can wipe out equity when leverage is high.
\end{prereq}

\subsection{Model Setup}

\textbf{Households} receive labor income and can invest in a risk-free asset or in intermediary equity. They cannot directly hold risky assets---they must invest through intermediaries.

\textbf{Intermediaries} hold risky assets (equities, bonds, derivatives) funded by household deposits (debt) and their own equity capital. They face an \textbf{equity capital constraint}: their equity must remain positive, and a tighter version requires equity to exceed a fraction of assets.

The capital ratio---the intermediary's equity divided by total assets---is the state variable:
\begin{equation}
w_t = \frac{E_t}{A_t}
\label{eq:capital-ratio-hk}
\end{equation}
\begin{itemize}[nosep]
  \item $w_t$ --- the intermediary's capital ratio at time $t$.
  \item $E_t$ --- equity capital (the intermediary's ``skin in the game'').
  \item $A_t$ --- total assets under management.
\end{itemize}

\subsection{Two Regimes}

The model generates two regimes, determined by how close $w_t$ is to the constraint boundary.

\textbf{When capital is ample} ($w_t$ far from the constraint): the intermediary is unconstrained. It can absorb shocks freely, and its SDF approximately equals the household SDF. Risk premia are moderate and stable.

\textbf{When capital is scarce} ($w_t$ near the constraint boundary): the intermediary's marginal value of wealth spikes. It must either sell risky assets (pushing prices down) or raise costly equity. Risk premia rise sharply. The SDF becomes:
\begin{equation}
\SDF_{t+1}^{\text{int}} \approx \SDF_{t+1}^{\text{HH}} + \mu_t \cdot \mathbf{1}\{\text{constraint near binding}\}
\label{eq:hk-sdf}
\end{equation}
where $\mu_t > 0$ is the shadow cost of the capital constraint (a Lagrange multiplier). The closer the intermediary is to the constraint, the larger $\mu_t$, and the more the intermediary SDF exceeds the household SDF.

\begin{prereq}[Lagrange Multipliers]
In constrained optimization, a \textbf{Lagrange multiplier} $\mu$ measures the ``shadow price'' of a constraint: how much the objective improves if you relax the constraint by one unit. If $\mu = 0$, the constraint is slack (non-binding) and irrelevant. If $\mu > 0$, the constraint is active and costly. In He-Krishnamurthy, $\mu_t$ is the shadow cost of the capital constraint. When $\mu_t = 0$, the intermediary is unconstrained and prices assets like a household. When $\mu_t > 0$, the constraint binds and the intermediary demands higher risk premia to compensate for its constrained state.
\end{prereq}

The nonlinearity is critical. Risk premia are \emph{not} a smooth, linear function of the capital ratio. Instead, they are nearly flat when capital is abundant and then rise steeply---almost like a hockey stick---as the constraint boundary approaches. This matches the empirical observation that risk premia are roughly stable during normal periods but spike violently during crises.

\begin{workedexample}{The Capital Constraint in Action}
Consider a simplified version of the He-Krishnamurthy setup. A dealer has \$10 billion in equity and \$90 billion in debt, for \$100 billion in total assets and a capital ratio of $w = 10\%$. Suppose the VaR limit is set so that the dealer must maintain $w \geq 5\%$.

\textbf{Normal times}: Assets drop 2\% (\$2B). Equity falls from \$10B to \$8B. Capital ratio: $8/98 = 8.2\%$---still well above the 5\% floor. The dealer can absorb the loss and continue trading. Risk premia change little.

\textbf{After a prior crisis}: Equity has already fallen to \$6B (capital ratio $= 6/96 = 6.25\%$, near the floor). The same 2\% asset decline now pushes equity to \$4.08B and the capital ratio to $4.08/94.08 = 4.3\%$---below the constraint. The dealer \emph{must} sell \$20B+ in risky assets to restore the ratio, flooding the market with supply. These forced sales push prices down further, creating a feedback loop. Risk premia spike.

\textbf{Key insight}: The same shock (2\% asset decline) produces a radically different risk-premium response depending on the starting capital ratio. This is the nonlinearity that He-Krishnamurthy's model captures, and it is what your VaR utilization feature is designed to measure.
\end{workedexample}

\begin{keyresult}[He and Krishnamurthy (2013)]
The model produces three predictions testable with SecDB data:
\begin{enumerate}[nosep]
  \item Risk premia are a \textbf{nonlinear} function of the intermediary capital ratio: low and stable when capital is ample, sharply elevated when capital is scarce.
  \item The speed of \textbf{reversion} from crisis to normal risk premia depends on the rate at which intermediaries rebuild capital (retained earnings, new equity issuance).
  \item \textbf{Policy implications}: injecting equity capital into intermediaries is more effective than reducing borrowing costs or purchasing distressed assets, because it directly relaxes the binding constraint.
\end{enumerate}
The calibrated model matches both the nonlinearity of risk premia during the 2007--2009 crisis and the speed of reversion back to normal levels.
\end{keyresult}

\begin{projectconnection}[VaR Utilization as the Capital Constraint Proxy]
He and Krishnamurthy's capital ratio $w_t$ is exactly what your VaR utilization feature measures. VaR utilization $=$ current VaR / VaR limit. When utilization approaches 100\%, the desk is at its constraint boundary---the model predicts risk premia should spike. You observe this daily from SecDB; the academics could only calibrate to aggregate quarterly data.
\end{projectconnection}

% ──────────────────────────────────────────────────────────────
\section{Adrian, Etula, and Muir (2014): Empirical Evidence}
\label{sec:adrian-etula-muir}
% ──────────────────────────────────────────────────────────────

He and Krishnamurthy (2013) is a theoretical model. Adrian, Etula, and Muir (2014, JF) take it to the data with a remarkable result.

\begin{prereq}[Cross-Sectional Pricing Test]
A cross-sectional pricing test asks: can a proposed factor explain \emph{why} different assets earn different average returns? You sort assets into portfolios, compute each portfolio's exposure (beta) to the factor, and check whether higher-beta portfolios earn higher average returns. The cross-sectional $R^2$ measures how much of the variation in average returns across portfolios is explained by the factor. Chapter~\ref{ch:asset-pricing} covered the Fama-MacBeth procedure for this.
\end{prereq}

\subsection{The Leverage Factor}

Adrian, Etula, and Muir construct a single factor: the \textbf{growth rate of broker-dealer leverage}. Their leverage measure comes from the Federal Reserve's Flow of Funds (Z.1) data:
\begin{equation}
\text{Leverage}_t = \frac{\text{Total financial assets}_t}{\text{Total financial assets}_t - \text{Total liabilities}_t}
\label{eq:aem-leverage}
\end{equation}
The numerator is total assets; the denominator is equity (assets minus liabilities). The factor is the innovation (unexpected change) in leverage growth, estimated quarterly from 1968Q1 to 2009Q4.

The logic: when funding conditions deteriorate, broker-dealers deleverage (sell risky assets, shrink balance sheets). This deleveraging \emph{is} the realization of the intermediary SDF being high. Assets that lose value precisely when dealers are deleveraging---i.e., assets with positive beta to leverage growth---are risky and should earn a premium.

\subsection{Headline Result}

\begin{keyresult}[Adrian, Etula, and Muir (2014)]
A single factor---broker-dealer leverage growth---prices 25 size and book-to-market equity portfolios, 10 momentum portfolios, and 6 Treasury bond portfolios (41 test assets total) with a cross-sectional $R^2$ of 77\% and an average annual pricing error of approximately 1\%. For comparison on the same test assets: the CAPM achieves $R^2 \approx 24\%$, and the Fama-French three-factor model achieves $R^2 \approx 73\%$. In the pre-crisis subsample (1968--2005), the intermediary factor still achieves $R^2 = 67\%$ with pricing errors of 37 basis points per year.
\begin{center}
\begin{tabular}{lccc}
\toprule
\textbf{Model} & \textbf{Factors} & \textbf{Cross-sectional $R^2$} & \textbf{Avg.\ pricing error} \\
\midrule
CAPM               & 1 (market) & 24\% & --- \\
Fama-French 3      & 3 (MKT, SMB, HML) & 73\% & --- \\
Intermediary (AEM) & 1 (leverage growth) & 77\% & ${\sim}1\%$/yr \\
\bottomrule
\end{tabular}
\end{center}
Sample: 1968Q1--2009Q4, quarterly. Data source: Fed Flow of Funds (Z.1), security broker-dealer sector.
\end{keyresult}

This is a striking result. One factor derived from broker-dealer balance sheets outperforms a three-factor model that was specifically \emph{designed} to price these portfolios. It strongly suggests that intermediary constraints, not household consumption, are what drives the cross-section of risk premia.

\begin{warning}[Data Limitations of the Leverage Factor]
The Fed Z.1 data is quarterly with a roughly two-month publication lag. This means the leverage factor is measured at a very low frequency and with stale information. For a trading signal, quarterly data with a two-month lag is nearly useless---by the time you observe the leverage change, markets have already moved. This is precisely the bottleneck your project solves: SecDB gives you the \emph{same economic information} (dealer balance-sheet constraints) at daily frequency with no publication lag.
\end{warning}

% ──────────────────────────────────────────────────────────────
\section{He, Kelly, and Manela (2017): Cross-Asset Evidence}
\label{sec:he-kelly-manela}
% ──────────────────────────────────────────────────────────────

Adrian, Etula, and Muir (2014) tested their factor on equities and Treasuries. He, Kelly, and Manela (2017, JFE) ask the harder question: does a single intermediary factor price \emph{all} asset classes?

\begin{prereq}[Primary Dealers]
\textbf{Primary dealers} are the roughly 20 financial institutions authorized to trade directly with the Federal Reserve Bank of New York in open-market operations. They include firms like Goldman Sachs, J.P.~Morgan, and Barclays. They are the backbone of U.S.\ fixed-income markets, intermediating trillions of dollars in Treasury, agency, and corporate bonds. Because they span multiple asset classes, their aggregate capital position reflects economy-wide intermediation capacity.
\end{prereq}

\subsection{The Capital Ratio Factor}

He, Kelly, and Manela construct an alternative to Adrian, Etula, and Muir's leverage measure. Instead of aggregate Flow of Funds leverage, they use the \textbf{equity capital ratio} of primary dealer holding companies:
\begin{equation}
\text{Capital ratio}_t = \frac{\text{Market equity}_t}{\text{Market equity}_t + \text{Book debt}_t}
\label{eq:hkm-capital-ratio}
\end{equation}
\begin{itemize}[nosep]
  \item Market equity is the stock-market capitalization of the primary dealer holding companies (from CRSP/Compustat).
  \item Book debt is total liabilities from financial statements.
  \item The factor is the \emph{innovation} (shock) to this ratio, estimated as the residual from an AR(1) model.
  \item Sample: 1970Q1--2012Q4, quarterly.
\end{itemize}

\begin{prereq}[AR(1) Models and Innovations]
An \textbf{AR(1)} (first-order autoregressive) model predicts a variable from its own lagged value: $x_t = \phi \, x_{t-1} + \varepsilon_t$, where $\phi$ is the persistence parameter and $\varepsilon_t$ is the residual (``innovation'' or ``shock''). The innovation $\varepsilon_t = x_t - \phi \, x_{t-1}$ is the \emph{unexpected} component of $x_t$---the part not predicted by the previous value. In asset pricing, innovations matter more than levels because markets anticipate persistent trends. A capital ratio of 12\% is not informative on its own; a \emph{surprise drop} from the expected 12\% to 10\% is what moves risk premia.
\end{prereq}

\begin{prereq}[Procyclical vs.\ Countercyclical]
A variable is \textbf{procyclical} if it rises during economic expansions and falls during contractions. GDP, employment, and stock prices are procyclical. A variable is \textbf{countercyclical} if it moves opposite to the cycle: unemployment, credit spreads, and the VIX are countercyclical. In intermediary asset pricing, the capital ratio is procyclical (rises in booms), while leverage $= 1/\text{capital ratio}$ is countercyclical (rises in crises). The He-Krishnamurthy model predicts this: dealers expand their balance sheets in good times and are forced to contract in bad times.
\end{prereq}

The capital ratio is procyclical: it rises in booms (equity values increase) and falls in crises (equity values collapse). This implies that intermediary \emph{leverage} is countercyclical---exactly as the He-Krishnamurthy model predicts.

\subsection{Seven Asset Classes}

\begin{keyresult}[He, Kelly, and Manela (2017)]
The capital ratio of primary dealers prices seven asset classes simultaneously:
\begin{enumerate}[nosep]
  \item U.S.\ equities (size, value, momentum portfolios)
  \item U.S.\ government bonds (maturity-sorted)
  \item U.S.\ corporate bonds (rating-sorted)
  \item Foreign sovereign bonds
  \item Options (volatility-sorted)
  \item Credit default swaps (CDS)
  \item Commodities and foreign exchange (FX)
\end{enumerate}
The capital risk premium---the compensation for bearing intermediary capital risk---ranges from 7\% to 11\% per year across five of the seven asset classes, with a pooled estimate of approximately 9\% per year. The price of risk is consistently positive and of similar magnitude when estimated separately for each asset class. Results are qualitatively stable in the pre-crisis subsample (1970Q1--2006Q4) and at monthly frequency.
\end{keyresult}

The fact that a \emph{single} intermediary factor works across equities, bonds, credit, FX, commodities, options, and sovereign debt is powerful evidence that intermediaries are the marginal investors in all these markets. This directly motivates your project's cross-asset panel design (Chapter~\ref{ch:panel}).

\begin{keyidea}[Cross-Asset Pricing Implies Cross-Asset Prediction]
If one intermediary factor prices seven asset classes simultaneously, then a shock to intermediary capital should predict returns across \emph{all} those asset classes---not just the one where the shock originated. Concretely: if the rates desk hits its VaR limit and is forced to sell Treasuries, the resulting capital depletion should also affect expected returns in credit, FX, and equities---because the same firm's capital backs all of its trading desks. This cross-asset spillover is what your ``cross-asset component VaR flow'' feature (Chapter~\ref{ch:features}) is designed to capture: you measure how VaR contributions shift across asset classes within the firm, and test whether those shifts predict returns in the receiving asset class.
\end{keyidea}

\subsection{AEM vs.\ HKM: Two Intermediary Measures}

Adrian, Etula, and Muir (AEM) and He, Kelly, and Manela (HKM) measure intermediary health differently:

\begin{center}
\begin{tabular}{p{3.5cm} p{5cm} p{5cm}}
\toprule
& \textbf{AEM (2014)} & \textbf{HKM (2017)} \\
\midrule
\textbf{Measure} & Book leverage of broker-dealer sector & Market equity capital ratio of primary dealers \\
\textbf{Data source} & Fed Z.1 Flow of Funds & CRSP/Compustat + Datastream \\
\textbf{Level/growth} & Leverage \emph{growth} innovations & Capital ratio \emph{level} innovations \\
\textbf{Cyclicality} & Procyclical leverage (rises in booms) & Procyclical capital ratio (rises in booms) \\
\textbf{Asset classes} & Equities + Treasuries (41 portfolios) & 7 asset classes \\
\textbf{Frequency} & Quarterly & Quarterly (monthly robustness) \\
\bottomrule
\end{tabular}
\end{center}

The two measures are conceptually related---both capture intermediary balance-sheet health---but they move differently in the data and can even give contradictory signals in some periods. He, Kelly, and Manela (2017) provide evidence that their capital ratio dominates AEM's leverage growth in joint cross-sectional tests across the broader set of asset classes.

\begin{warning}[The AEM vs.\ HKM Debate Is Unresolved]
The academic literature has not settled whether leverage growth (AEM) or the capital ratio level (HKM) is the ``correct'' intermediary state variable. The two can give different signs: AEM leverage is procyclical (leverage rises in booms), while HKM's capital ratio is also procyclical (it rises in booms)---but the innovations (shocks) move in opposite directions for some sample periods. For your project, this debate matters less than the shared conclusion: \emph{some} measure of dealer balance-sheet health is a priced state variable. SecDB gives you multiple such measures simultaneously (VaR utilization, component VaR, factor concentration), so you can test which one carries the most information in your sample without taking sides in the AEM-vs.-HKM debate.
\end{warning}

% ──────────────────────────────────────────────────────────────
\section{Adrian and Shin (2010): Repos and VIX}
\label{sec:adrian-shin}
% ──────────────────────────────────────────────────────────────

While He-Krishnamurthy and AEM/HKM focus on the cross-section of expected returns, Adrian and Shin (2010, JFI) provide direct evidence on the \emph{time-series} forecasting power of dealer balance sheets.

\begin{prereq}[Repurchase Agreements (Repos)]
A \textbf{repo} (repurchase agreement) is a short-term collateralized loan. A dealer sells a security (e.g., a Treasury bond) to a counterparty and agrees to buy it back the next day (or within a few days) at a slightly higher price. The price difference is the interest rate. Repos are the primary funding mechanism for dealer inventories: a dealer with \$100 billion in bond inventory might fund \$95 billion of it through overnight repos. Changes in total repo outstanding directly reflect changes in dealer balance-sheet size---when repos grow, dealers are expanding; when repos shrink, dealers are deleveraging.
\end{prereq}

\begin{prereq}[The VIX Index]
The \textbf{VIX} (Chicago Board Options Exchange Volatility Index) measures the market's expectation of 30-day S\&P~500 volatility, implied from option prices. It is often called the ``fear index.'' A VIX of 15 means the market expects annualized volatility of 15\%. The long-run average is roughly 19--20. During crises, VIX spikes: it reached 80 in October 2008 and 66 in March 2020. A \textbf{VIX innovation} is the unexpected change in VIX---the component not predicted by its own past values.
\end{prereq}

Adrian and Shin study whether changes in dealer repo positions forecast changes in market risk. Their logic: when dealers expand their balance sheets via repos, they are absorbing more risk. When they contract, they are pulling back. The aggregate repo position is a real-time measure of dealer risk appetite.

\begin{keyresult}[Adrian and Shin (2010)]
Growth in broker-dealer repo positions forecasts VIX innovations. Using weekly data from January 1990 to August 2007, a regression of VIX changes on lagged repo growth yields an $R^2$ of 11.6\%, compared to 8.9\% using only the lagged VIX level as a predictor. The incremental $R^2$ from repo growth---2.7 percentage points---may seem small, but at weekly frequency for a financial variable, this is economically meaningful. The result is stronger for the \textbf{volatility risk premium} (the gap between implied and realized volatility), which has a direct interpretation as the price of risk.
\begin{center}
\begin{tabular}{lcc}
\toprule
\textbf{Predictor} & \textbf{$R^2$ (VIX changes)} & \textbf{Interpretation} \\
\midrule
Lagged VIX level only & 8.9\% & Mean reversion in VIX \\
+ Repo growth & 11.6\% & Dealer balance sheets add info \\
Incremental $R^2$ & +2.7 pp & Dealer risk appetite signal \\
\bottomrule
\end{tabular}
\end{center}
Data: weekly, January 1990--August 2007, U.S.\ primary dealers.
\end{keyresult}

\begin{intuition}[Why Repos Track Dealer Risk Appetite]
Imagine you run a bond desk. To hold \$1 billion in Treasuries, you need funding. You repo out the bonds overnight---lending them as collateral and receiving \$990 million in cash. Tomorrow, you buy them back for \$990.01 million (the overnight interest). If you want to hold \emph{more} bonds, you do more repos. If you want to reduce exposure, you unwind repos. At the aggregate level, total repos outstanding is therefore a near-real-time measure of how much balance sheet the dealer sector is deploying. When repos grow, dealers are leaning into risk. When repos shrink, dealers are pulling back. Adrian and Shin's result says that this risk-appetite measure forecasts future volatility---because when dealers are aggressively expanding, they are absorbing risks that may later unwind.
\end{intuition}

This result is directly relevant to your project for two reasons:

\begin{enumerate}[nosep]
  \item It establishes that \textbf{VIX innovations} are the cleanest target variable for testing intermediary signals. The VIX represents the market price of volatility risk, and dealer balance sheets predict it. Your Phase~2 feature testing will use VIX innovations as one of four targets (Section~\ref{sec:project-targets}).
  \item It provides a \textbf{benchmark} for your features. Adrian and Shin achieve 2.7pp of incremental $R^2$ with quarterly-lagged, aggregate repo data. With daily, granular VaR and scenario data from SecDB, your features should contain at least as much information---and likely more. If your features cannot beat 2.7pp incremental $R^2$ for VIX forecasting, something is wrong with your feature engineering, not with the theory.
\end{enumerate}

% ──────────────────────────────────────────────────────────────
\section{Adrian and Brunnermeier (2016): CoVaR}
\label{sec:covar}
% ──────────────────────────────────────────────────────────────

\begin{prereq}[Value at Risk (VaR)]
\textbf{Value at Risk} at confidence level $\alpha$ is the loss threshold such that the probability of exceeding it is $\alpha$. Formally, $\VaR_\alpha(X) = -\inf\{x : P(X \leq x) > \alpha\}$ for a return random variable $X$. In plain language: the 99\% daily VaR of \$10 million means ``we expect to lose more than \$10 million on only 1\% of trading days.'' Chapter~\ref{ch:risk-systems} covers VaR computation in detail; here you only need the concept.
\end{prereq}

Adrian and Brunnermeier (2016, AER) introduce \textbf{CoVaR}---the VaR of the financial system \emph{conditional on} a particular institution being in distress. This shifts VaR from a firm-level risk measure to a \emph{systemic} risk measure.

\begin{definition}[CoVaR and $\Delta$CoVaR]
The \textbf{CoVaR} of the financial system conditional on institution $j$ is:
\begin{equation}
P\!\left(R_{\text{system}} \leq \text{CoVaR}_{q}^{\text{system}|j} \;\middle|\; R_j = \VaR_q^j\right) = q
\label{eq:covar-def}
\end{equation}
where $R_{\text{system}}$ is the return on the financial system and $R_j$ is institution $j$'s return.
\begin{itemize}[nosep]
  \item $\text{CoVaR}_{q}^{\text{system}|j}$ --- the $q$-quantile loss of the system, given that institution $j$ is at its own VaR (i.e., in distress).
  \item $q$ --- the tail quantile (typically $q = 0.05$ or $0.01$).
\end{itemize}

The \textbf{$\Delta$CoVaR} of institution $j$ is the difference between the system's CoVaR when $j$ is in distress versus when $j$ is at its median state:
\begin{equation}
\Delta\text{CoVaR}_{q}^{j} = \text{CoVaR}_{q}^{\text{system}|j \text{ at VaR}} - \text{CoVaR}_{q}^{\text{system}|j \text{ at median}}
\label{eq:delta-covar}
\end{equation}
$\Delta$CoVaR measures institution $j$'s \emph{contribution} to systemic risk: how much worse the system's tail risk becomes when $j$ is in trouble.
\end{definition}

\subsection{Key Findings}

\begin{keyresult}[Adrian and Brunnermeier (2016)]
Four institution-level characteristics significantly predict $\Delta$CoVaR:
\begin{enumerate}[nosep]
  \item \textbf{Leverage} --- higher leverage $\Rightarrow$ larger systemic risk contribution.
  \item \textbf{Size} --- larger institutions contribute more to systemic risk.
  \item \textbf{Maturity mismatch} --- funding long-term assets with short-term debt increases fragility.
  \item \textbf{Asset price booms} --- rapid growth in asset values signals building vulnerability.
\end{enumerate}
Critically, $\Delta$CoVaR is \textbf{countercyclical}: it builds up during booms (when measured VaR is \emph{low}) and realizes during crises. A forward-looking version of $\Delta$CoVaR measured in 2006Q4 would have predicted more than one-third of the realized systemic losses during the 2007--2009 financial crisis.
\end{keyresult}

\subsection{Why CoVaR Matters for Your Features}

\begin{projectconnection}[CoVaR and Your VaR Features]
Your SecDB VaR features capture a version of the CoVaR logic at the desk level:
\begin{itemize}[nosep]
  \item \textbf{Firm-level delta VaR} (how the desk's trades change total firm VaR) is analogous to $\Delta$CoVaR---it measures one desk's contribution to the firm's aggregate tail risk.
  \item \textbf{VaR utilization} (current VaR / VaR limit) captures the leverage-like constraint channel that Adrian and Brunnermeier identify as the strongest predictor.
  \item The countercyclical property of $\Delta$CoVaR matches the He-Krishnamurthy prediction: risk builds when constraints are slack and realizes when they bind.
\end{itemize}
You will not compute CoVaR directly (it requires cross-institution data), but your features capture the same underlying economics at higher frequency and finer granularity.
\end{projectconnection}

% ──────────────────────────────────────────────────────────────
\section{Coval and Stafford (2007): Fire Sales and Crowding}
\label{sec:coval-stafford}
% ──────────────────────────────────────────────────────────────

The intermediary theories above focus on aggregate dealer capital. Coval and Stafford (2007, JFE) provide the microeconomic mechanism for \emph{how} constrained intermediaries affect prices: forced selling.

\begin{prereq}[Mutual Fund Flows]
When investors withdraw money from a mutual fund (``outflows''), the fund must sell holdings to raise cash. When many investors withdraw simultaneously, the fund is a \textbf{forced seller}---it must sell regardless of price. This is distinct from an informed seller who sells because they believe the asset is overvalued. Forced selling creates temporary price pressure that reverses once the selling stops.
\end{prereq}

Coval and Stafford study mutual funds experiencing extreme outflows. When funds are forced to liquidate positions, they sell stocks mechanically. The key finding:

\begin{keyresult}[Coval and Stafford (2007)]
Stocks heavily sold by distressed mutual funds (those experiencing outflows in the bottom 5\% of the distribution) experience significant price declines that subsequently reverse. The fire-sale discount is approximately 10--15\% over the selling period, with roughly half reversing within 6--12 months. The price impact is concentrated in stocks where the distressed fund holds a large share of daily volume---i.e., where the forced selling overwhelms the market's absorptive capacity.
\end{keyresult}

The connection to intermediary asset pricing is direct: forced selling is the mechanism through which capital constraints transmit to prices. When a dealer hits its VaR limit, it must reduce risk---often by selling positions. If many dealers are constrained simultaneously (as in 2008), the forced selling is correlated across firms, amplifying the price impact.

\begin{projectconnection}[Factor Concentration and Fire-Sale Risk]
Your \textbf{factor concentration} feature (the Herfindahl index of factor-VaR contributions) captures fire-sale vulnerability directly. When factor concentration is high---meaning most of the desk's VaR comes from a few correlated positions---a constraint-driven forced sale will hit the same positions and the same markets. Low concentration means the desk can reduce risk by trimming many small, diverse positions with minimal market impact. High Herfindahl $\Rightarrow$ high fire-sale risk $\Rightarrow$ higher expected risk premia. This is a testable prediction from Coval-Stafford combined with He-Krishnamurthy.
\end{projectconnection}

% ──────────────────────────────────────────────────────────────
\section{The Data Bottleneck---Why Your Project Has an Edge}
\label{sec:data-bottleneck}
% ──────────────────────────────────────────────────────────────

Every paper in this chapter shares the same limitation: they use public, low-frequency, backward-looking data to measure intermediary constraints.

\begin{center}
\begin{tabular}{p{3.5cm} p{2.5cm} p{2cm} p{2cm} p{3cm}}
\toprule
\textbf{Paper} & \textbf{Data Source} & \textbf{Freq.} & \textbf{Lag} & \textbf{Granularity} \\
\midrule
He-Krishnamurthy (2013) & Calibration only & --- & --- & Aggregate \\
Adrian-Etula-Muir (2014) & Fed Z.1 & Quarterly & ${\sim}2$ months & Sector aggregate \\
He-Kelly-Manela (2017) & CRSP/Compustat & Quarterly & ${\sim}1$ month & Firm (holding co.) \\
Adrian-Shin (2010) & Fed Z.1 & Weekly & ${\sim}1$ week & Sector aggregate \\
Adrian-Brunnermeier (2016) & CRSP, Call Reports & Quarterly & ${\sim}2$ months & Firm \\
\midrule
\textbf{Your project} & \textbf{SecDB} & \textbf{Daily} & \textbf{T+0} & \textbf{Desk / position} \\
\bottomrule
\end{tabular}
\end{center}

The advantage is threefold:

\begin{enumerate}[nosep]
  \item \textbf{Frequency}: Daily vs.\ quarterly gives you roughly 60$\times$ more observations per year. For ML methods that need data, this matters.
  \item \textbf{Timeliness}: You observe today's risk state today. The academics observe last quarter's risk state two months later. Any predictive signal in intermediary constraints has likely dissipated by the time public data is released.
  \item \textbf{Granularity}: You observe risk decomposed by desk, factor, scenario, and position. The academics observe a single aggregate number for the entire broker-dealer sector. Disaggregation lets you identify \emph{which} constraints matter and \emph{where} the pressure is building.
\end{enumerate}

\begin{warning}[Better Data Does Not Guarantee Better Results]
Three caveats temper the data advantage:
\begin{enumerate}[nosep]
  \item \textbf{Single firm}: You observe one dealer (Goldman Sachs), not the entire system. If the theory requires \emph{aggregate} intermediary constraints, one firm's data may be a noisy proxy.
  \item \textbf{Short history}: SecDB data may span only 5--10 years of daily observations, versus the 40+ year samples in the academic literature. Short samples increase overfitting risk and reduce the power of statistical tests. This is why the Deflated Sharpe Ratio (Chapter~\ref{ch:validation}) is non-negotiable.
  \item \textbf{Endogeneity}: The same firm whose risk data you observe is also a market participant. Its risk-taking affects prices, and prices affect its risk measures. Disentangling cause and effect is harder with internal data than with external data.
\end{enumerate}
Acknowledge these limitations in your pitch (Phase~0) and your final report (Phase~5). The correct framing is: ``better data that \emph{should} contain more signal, pending rigorous validation''---not ``we have the best data so the signal is guaranteed.''
\end{warning}

% ──────────────────────────────────────────────────────────────
\section{From Theory to Features: The Mapping}
\label{sec:theory-to-features}
% ──────────────────────────────────────────────────────────────

Let us make the connection between the theory in this chapter and the features you will build in Phase~2 (Chapter~\ref{ch:features}) fully explicit.

\begin{center}
\begin{tabular}{p{2.8cm} p{3cm} p{3cm} p{4.5cm}}
\toprule
\textbf{Theory Concept} & \textbf{Academic Proxy} & \textbf{Your SecDB Feature} & \textbf{Prediction} \\
\midrule
Capital constraint (He-Krishnamurthy) & Aggregate capital ratio & VaR utilization (\%) & High utilization $\Rightarrow$ higher risk premia \\
\addlinespace
Balance-sheet expansion (Adrian-Shin) & Repo growth & Delta VaR, component VaR & Rising VaR $\Rightarrow$ dealers absorbing risk $\Rightarrow$ future volatility \\
\addlinespace
Leverage factor (AEM) & Sector leverage growth & VaR rate-of-change & Rapid VaR changes $\Rightarrow$ leverage adjustment \\
\addlinespace
Cross-asset pricing (HKM) & Capital ratio across 7 classes & Cross-asset component VaR shifts & Risk reallocation between desks $\Rightarrow$ return predictability \\
\addlinespace
Crowding / fire sales (Coval-Stafford) & Mutual fund flow pressure & Factor concentration (HHI) & High Herfindahl $\Rightarrow$ crowded risk $\Rightarrow$ fragile \\
\addlinespace
Systemic risk (Adrian-Brunnermeier) & $\Delta$CoVaR & Firm-level delta VaR & Desk contribution to firm risk \\
\addlinespace
Tail risk (He-Krishnamurthy) & Crisis vs.\ normal calibration & Scenario P\&L dispersion & Wide scenario spread $\Rightarrow$ tail uncertainty \\
\bottomrule
\end{tabular}
\end{center}

Each row is a testable hypothesis. In Phase~2, you will test each one against four targets: VIX innovations, asset-class drawdowns, realized volatility, and cross-asset momentum reversals. The theory in this chapter tells you \emph{why} each feature should predict---without that ``why,'' you are data-mining.

\label{sec:project-targets}

\begin{keyidea}[Theory-Motivated Features vs.\ Data-Mining]
The most important methodological principle in this project is: \textbf{every feature must be motivated by theory before you look at the data}. This chapter provides that motivation. When you present results in Phase~5, the narrative is not ``my ML model found a pattern.'' The narrative is: ``Intermediary asset pricing theory (He-Krishnamurthy 2013, AEM 2014, HKM 2017) predicts that dealer constraints drive risk premia. I tested this prediction using high-frequency proprietary data and found [result].'' The first framing invites the response ``that is probably overfit.'' The second invites ``tell me more.''
\end{keyidea}

% ──────────────────────────────────────────────────────────────
\section{Summary}
\label{sec:ch2-summary}
% ──────────────────────────────────────────────────────────────

\begin{enumerate}[nosep]
  \item The \textbf{equity premium puzzle} shows that household consumption is too smooth to explain observed risk premia.
  \item \textbf{Intermediary asset pricing} resolves this by making the financial intermediary---not the household---the marginal investor whose constraints drive the SDF.
  \item \textbf{He and Krishnamurthy (2013)} provide the theoretical model: the intermediary's capital ratio is the state variable; risk premia spike when capital is scarce and the equity constraint binds.
  \item \textbf{Adrian, Etula, and Muir (2014)} show empirically that a single broker-dealer leverage factor prices equities, momentum, and bonds with $R^2 = 77\%$, outperforming the CAPM ($24\%$) and matching Fama-French ($73\%$). Sample: 1968Q1--2009Q4, quarterly.
  \item \textbf{He, Kelly, and Manela (2017)} extend the result across seven asset classes (equities, government and corporate bonds, sovereign bonds, options, CDS, commodities, FX) using the primary-dealer capital ratio. The capital risk premium is approximately 9\% per year. Sample: 1970Q1--2012Q4.
  \item \textbf{Adrian and Shin (2010)} demonstrate that repo growth forecasts VIX innovations ($R^2$ rises from 8.9\% to 11.6\% when repos are included). This motivates VIX innovations as your cleanest target variable.
  \item \textbf{Adrian and Brunnermeier (2016)} define CoVaR---systemic VaR---and show that leverage, size, and maturity mismatch predict systemic risk contributions. A forward-looking measure from 2006Q4 would have predicted more than one-third of realized crisis losses.
  \item All academic work uses \textbf{quarterly, lagged, aggregate} data. Your project uses \textbf{daily, real-time, desk-level} SecDB data---a frequency and granularity advantage of roughly 60$\times$.
  \item \textbf{Coval and Stafford (2007)} show that forced selling by constrained funds creates 10--15\% fire-sale discounts that partially reverse. Factor concentration (Herfindahl) measures vulnerability to this mechanism.
  \item \textbf{Better data alone is not enough.} Single-firm data, short history, and endogeneity are real limitations.
  \item Every Phase~2 feature maps to a specific theoretical prediction from this chapter. Theory-motivated features are your defense against overfitting and your strongest pitch to the desk.
  \item \textbf{Publication decay caveat}: McLean and Pontiff (2016) show that published anomalies lose 30--50\% of their magnitude post-publication. The intermediary factors here may have partially decayed. Your project's edge is not the theory (which is public) but the proprietary data that measures it at higher resolution.
\end{enumerate}

% ──────────────────────────────────────────────────────────────
\section*{Key Results Reference}
\label{sec:ch2-key-results}
% ──────────────────────────────────────────────────────────────

\begin{center}
\begin{tabular}{p{3.5cm} p{3.5cm} p{6.5cm}}
\toprule
\textbf{Paper} & \textbf{Key Claim} & \textbf{Headline Result} \\
\midrule
He \& Krishnamurthy (2013, AER) & Intermediary capital ratio drives risk premia & Theoretical model: risk premia nonlinear in capital ratio; spike when constraint binds; calibrated to crisis dynamics \\
\addlinespace
Adrian, Etula \& Muir (2014, JF) & Leverage growth is a priced SDF factor & Single factor: cross-sectional $R^2 = 77\%$ on 41 portfolios; avg.\ pricing error ${\sim}1\%$/yr; CAPM $R^2 = 24\%$, FF3 $R^2 = 73\%$ \\
\addlinespace
He, Kelly \& Manela (2017, JFE) & One factor prices all asset classes & Capital ratio prices 7 asset classes; capital risk premium ${\approx}9\%$/yr; positive price of risk in each class separately \\
\addlinespace
Adrian \& Shin (2010, JFI) & Repos forecast VIX & Repo growth raises VIX forecasting $R^2$ from 8.9\% to 11.6\% (weekly, 1990--2007) \\
\addlinespace
Adrian \& Brunnermeier (2016, AER) & CoVaR measures systemic risk & Leverage, size, maturity mismatch predict $\Delta$CoVaR; 2006Q4 measure predicts ${>}\frac{1}{3}$ of 2007--09 realized losses \\
\addlinespace
Coval \& Stafford (2007, JFE) & Forced selling creates fire-sale discounts & 10--15\% discount on stocks sold by distressed funds; ${\sim}50\%$ reversal within 6--12 months \\
\bottomrule
\end{tabular}
\end{center}
